% Generated draft: requires mathematical and source review.
\documentclass[11pt]{article}
\usepackage[margin=25mm]{geometry}
\usepackage{fontspec}
\setmainfont{Latin Modern Roman}
\usepackage{amsmath,amssymb,mathtools,unicode-math}
\setmathfont{Latin Modern Math}
\usepackage{enumitem,needspace,xurl,hyperref}
\setlength{\parindent}{0pt}
\setlength{\parskip}{5pt}
\setlength{\emergencystretch}{3em}
\newcommand{\N}{\mathbb{N}}
\newcommand{\Z}{\mathbb{Z}}
\newcommand{\Q}{\mathbb{Q}}
\newcommand{\R}{\mathbb{R}}
\newcommand{\C}{\mathbb{C}}
\newcommand{\F}{\mathbb{F}}
\newcommand{\A}{\mathbb{A}}
\newcommand{\PP}{\mathbb{P}}
\newcommand{\E}{\mathbb{E}}
\newcommand{\eps}{\varepsilon}
\newcommand{\dd}{\,\mathrm{d}}
\begin{document}
\section*{UnsolvedMath ID: 30001217}
\textbf{Source problem code:} catalogue entry
\subsection{Definitions and mathematical statement}
Let \(\mathfrak g\) be a simple Lie algebra of type \(A\), \(D\), or \(E\). Write \(\mathcal C_\ell\) for the monoidal subcategory of finite-dimensional quantum-affine \(\mathfrak g\)-modules intended by the source, and write \(R_\ell=K_0(\mathcal C_\ell)\) for its Grothendieck ring. Here \([M]\) denotes the class of an object \(M\), the additive relations are those arising from short exact sequences, and multiplication is induced by tensor product: \([M][N]=[M\otimes N]\).

For a specified cluster algebra structure on \(R_\ell\), let \(\operatorname{CM}(R_\ell)\) denote its cluster monomials: products of nonnegative integral powers of variables belonging to a single cluster, subject to the chosen convention concerning frozen variables. A nonzero module is simple if it has no nonzero proper submodule. A simple object \(S\) is called real if \(S\otimes S\) is simple.

The source asks whether, for every permitted \(\mathfrak g\) and level \(\ell\), both assertions hold:
\[
\forall m\in\operatorname{CM}(R_\ell),\quad
\exists S\in\mathcal C_\ell\text{ simple such that }[S]=m,
\]
\[
\forall S\in\mathcal C_\ell\text{ simple},\quad
[S]\in\operatorname{CM}(R_\ell)
\quad\Longleftrightarrow\quad
S\otimes S\text{ is simple}.
\]
The supplied record does not specify the ground field, quantum parameter, precise quantum-affine algebra and module convention, allowed levels, defining restrictions of \(\mathcal C_\ell\), or cluster algebra structure and coefficient convention on \(R_\ell\). Consequently these formulas express the requested assertions conditional on those source-dependent choices; they do not supply a complete independent specification of the objects being quantified over.
\subsection{Short English statement}
The question asks whether every cluster monomial is the class of a simple quantum-affine module in the intended category, and whether a simple module has a cluster-monomial class exactly when its tensor square is simple. The category, level range, and cluster structure require clarification from the cited sources. The supplied status is partially solved; this draft does not independently verify that assessment.
\subsection{Sources}
\begin{itemize}
\item Local UnsolvedMath registry: supplied complete source record for permanent ID \(30001217\), titled ``Cluster Monomials and Simple Quantum Affine Modules,'' in the Oberwolfach Reports Open Problems collection; exact source problem code \texttt{catalogue entry}. The recorded proposal year is 2009; no individual proposer is supplied.
\item Supplied report citation: \emph{Enveloping Algebras and Geometric Representation Theory} (2009), Oberwolfach Reports, pp.\ 801--866, DOI 10.4171/owr/2009/15. \url{https://doi.org/10.4171/owr/2009/15}. The supplied background describes these reports as records of mathematical workshops and their open-problem discussions.
\item Supplied literature reference, labeled ``Hernandez--Leclerc primary paper'': \url{https://arxiv.org/abs/0903.1452}.
\item Supplied literature reference, labeled ``Hernandez--Leclerc current survey'': \url{https://arxiv.org/abs/1902.01432}.
\item Supplied literature reference, labeled ``original OWR report'': \url{https://doi.org/10.4171/owr/2009/15}.
\item Literature assessment dated exactly \texttt{2026-08-21}: the supplied curation labels the problem partially solved and reports that the cluster-monomial/simple-object assertion is proved in important finite-type and selected quantum-affine cases, while primary surveys continue to describe the arbitrary-\(\mathfrak g\), arbitrary-level version as conjectural. This is supplied dated curation derived from the OWR record, not an independently verified assessment; its time-sensitive status and cited claims require independent verification.
\item The registry labels the statement \texttt{corrected\_verified}, explaining that the complete statement was preserved while malformed or undelimited notation was normalized. This supplied editorial label does not establish mathematical correctness or resolution.
\end{itemize}
\end{document}
